\documentclass[10pt,a4paper]{amsart}
\usepackage[margin=19mm]{geometry}
\usepackage{amsmath,amssymb,mathtools}
\usepackage[scr=rsfs,cal=euler]{mathalfa}
\usepackage{xfrac}
\usepackage[unicode,hidelinks,bookmarks=false]{hyperref}
\newcommand{\ie}{i.e.\ }
\newcommand{\cf}{cf.\ }
\newcommand{\resp}{respectively\ }
\newcommand{\Subsection}{\subsection}
\providecommand{\nobreakdash}{\nobreak\hbox{-}}
\setlength{\emergencystretch}{2em}
\multlinegap0pt
\usepackage{amssymb}
\usepackage{mathtools}
\usepackage[shortlabels]{enumitem}
\usepackage{xcolor}
\newtheorem{lemma}{Lemma}
\newtheorem{proposition}{Proposition}
\newcommand{\A}{\mathbb A}
\newcommand{\R}{\mathbb R}
\newcommand{\T}{\mathbb T}
\newcommand{\Z}{\mathbb Z}
\newcommand{\abs}[1]{\left\lvert #1\right\rvert}
\newcommand{\eps}{\varepsilon}
\newcommand{\la}{\lambda}
\newcommand{\norm}[1]{\left\lVert #1\right\rVert}
\newcommand{\supp}{\operatorname{supp}}
\title{Non-unique singular solutions for KP and modified KP equations on $\mathbb T^2$ and $\mathbb R^2$}
\author{Alexandru F. Radu}
\date{Source: arXiv:2608.20972v1 (CC BY 4.0)}
\begin{document}
\maketitle

\noindent\emph{Source and licence.} Non-unique singular solutions for KP and modified KP equations on $\mathbb T^2$ and $\mathbb R^2$. Version arXiv:2608.20972v1 (CC BY 4.0), distributed by arXiv under the Creative Commons Attribution 4.0 International licence, http://creativecommons.org/licenses/by/4.0/, as recorded in the arXiv licence metadata for this exact version and shown on the abstract page https://arxiv.org/abs/2608.20972v1. Full licence notice: \texttt{licenses/arxiv-2608.20972-CC-BY-4.0.txt}.\par\smallskip

\noindent\textit{Editorial scope.} Counted unit: the article's proposition prop:Euclidean-onestep (source0.tex lines 2169--2228). The article's complete original proof of it is delivered with it (source0.tex lines 2230--2495, one \texttt{proof} environment of 266 lines). No mathematics is written, repaired or re-derived: the statement and the proof are the article's own lines, in the article's own notation, and no retained line is substituted or re-flowed.  Retained uncounted as the source-local context the proof needs: lines 661--665; lines 676--680; lines 697--721; lines 875--880; lines 1482--1487; lines 1607--1616; lines 1626--1629; lines 1737--1740; lines 1747--1750; lines 1761--1763; lines 1766--1772; lines 1780--1786; lines 1788--1791; lines 1803--1823; lines 1809--1813; lines 1858--1863; lines 2004--2009; lines 2039--2085; lines 2089--2093.  Nothing was omitted from any retained span, so the package carries no omission record.  No retained line contains a citation, so the wrapper carries no bibliography and no bibliography entry.  No retained line contains a figure environment, an image-inclusion command or drawing code, so no image is delivered and the package declares no asset dependency.  Editorial formatting only, in addition: an AMS \texttt{amsart} preamble that restates the article's own declarations and macros used by the retained lines, namely the environments \texttt{lemma}, \texttt{proposition} on separate counters, together with the standard packages those lines need.  The retained environments print in this document as Lemma 1, Proposition 1 in order of appearance, whereas the article prints the same items with the numbering of its own full document.  The complete native source of the article as distributed in the arXiv e-print for this exact version is bundled unchanged as \texttt{sources/208215/source0.tex}.
\begin{equation}\label{eq:relaxed-general}
 \partial_tu+(-1)^{(d+1)/2}\partial_x^du
 +\kappa\partial_x^{-1}\partial_y^2u+\partial_x(u^r)
 =\partial_xE.
\end{equation}

\begin{equation}\label{eq:Euclidean-update-general}
 E^+=E+(u+w)^r-u^r+(-1)^{(d+1)/2}\partial_x^{d-1}w
 +\kappa\partial_x^{-2}\partial_y^2w
 +\partial_t\partial_x^{-1}w.
\end{equation}

\begin{lemma}\label{lem:slab}
Fix $0<\eps<1$.  For every sufficiently large
$\la>1$ such that $\mu=\la^\eps\in\mathbb N$, set
\begin{equation*}
 \nu=\la^{1+\eps}=\la\mu.
\end{equation*}
We call these choices of $\la$ admissible.
Then there exists a real, even trigonometric polynomial
$\rho_{\la,\eps}=\rho_{\la,\eps}(x)$ such that
\begin{align}
 &\int_\T\rho_{\la,\eps}\,dx=0,
 \qquad \int_\T\rho_{\la,\eps}^2\,dx=1,
 \label{eq:slabmoments}\\
 &\supp\widehat\rho_{\la,\eps}
 \subset(\mu\Z\setminus\{0\})\cap[-2\nu,2\nu],
 \notag\\
 &\norm{\rho_{\la,\eps}}_{L^p(\T)}
 \leq C_{p,\eps}\la^{(1-\eps)(1/2-1/p)},
 \qquad1\leq p\leq\infty,
 \label{eq:slabLp}\\
 &0\leq\widehat\rho_{\la,\eps}(n)
 \leq C_\eps\la^{-(1-\eps)/2}.
 \label{eq:slab-coefficient-bound}
\end{align}
\end{lemma}

\begin{equation}\label{eq:time-cutoff-lemma}
   0\leq h\leq1,
   \qquad h=1\ \text{on }I,
   \qquad \supp h\subset(t_--\tau,t_++\tau),
   \qquad \norm{h'}_{L^\infty}\leq C\tau^{-1},
\end{equation}

\begin{equation}\label{eq:Euclidean-periodic-decoupling}
    \norm{aU}_{L^p(\R^2)}
    \lesssim
    \norm{a}_{W^{3,p}(\R^2)}
    \norm{U}_{L^p(\T^2)}.
\end{equation}

\begin{equation}\label{eq:Euclidean-spatial-cutoff-moment}
\begin{gathered}
   \widehat{\chi_L}\geq0,
   \qquad
   \int_{\R^2}\widehat{\chi_L}(k)\,dk=1,\\
   \int_{\R^2}\abs z
      \bigl(\widehat{\chi_L}*\widehat{\chi_L}\bigr)(z)\,dz
      \leq\frac{C_\chi}{L}.
\end{gathered}
\end{equation}

\begin{equation}\label{eq:localized-constant-scaling}
    \norm{\partial_x(A\chi_L^2)}_{\A^{-S}_{\R^2}}
      \leq \frac{C_\chi A}{L}.
\end{equation}

\begin{equation}\label{eq:Euclidean-A-choice}
    A\geq C_M\left(1+\norm{E}_{C_tL^\infty}
                    +\norm{E}_{\A^M_{\R^2,t}}\right),
\end{equation}

\begin{equation}\label{eq:Euclidean-amplitude}
    a_L=\chi_L(A-E)^{1/2},
    \qquad a_{L,\ell}=P_{\leq\ell}^{x,y}a_L.
\end{equation}

\begin{equation}\label{eq:Euclidean-increment}
    w(t,x,y)=h(t)a_{L,\ell}(t,x,y)\rho_{\la,\eps}(x).
\end{equation}

\begin{equation}\label{eq:Euclidean-cone}
    \supp\widehat w(t)\subset
    \left\{(\xi,\eta):
       \mu-2\ell\leq\abs\xi\leq2\nu+2\ell,
       \quad \abs\eta\leq2\ell
    \right\}.
\end{equation}

\begin{equation}\label{eq:Euclidean-Lp}
    \norm{w}_{C_tL^p(\R^2)}
    \lesssim
    \norm{h a_{L,\ell}}_{C_tW^{3,p}(\R^2)}
    \la^{(1-\eps)(1/2-1/p)},
    \qquad 1\leq p<2.
\end{equation}

\begin{equation}\label{eq:Euclidean-cancellation}
 E+w^2=h^2\Bigl((1-\chi_L^2)E+A\chi_L^2
 +(a_{L,\ell}^2-a_L^2)+a_{L,\ell}^2(\rho_{\la,\eps}^2-1)\Bigr).
\end{equation}

\begin{proposition}
\label{prop:Euclidean-amplitude-exact-coefficient}
Fix $S\geq1$ and $M>S+6$.  Let $E$ be real with
$\norm{E}_{\A^M_{\R^2,t}}<\infty$, and let $0\leq h\leq1$ be a smooth
function of time satisfying $hE=E$.  Choose $A$ as in
\eqref{eq:Euclidean-A-choice}, with the constant there large enough that
\begin{equation}\label{eq:Euclidean-amplitude-smallness}
   A\geq2\norm E_{C_tL^\infty},
   \qquad
   \frac{C_{\mathrm{alg},M}\norm E_{\A^M_{\R^2,t}}}{A}\leq\frac12.
\end{equation}
For $L,\ell\geq1$, define $a_L$ and $a_{L,\ell}$ by
\eqref{eq:Euclidean-amplitude}.  Then
\begin{equation}\label{eq:Euclidean-amplitude-exact-coefficient}
   \operatorname{AF}^{\partial_x}_S(h a_{L,\ell})
   \leq
   \norm{\partial_xE}_{\A^{-S}_{\R^2,t}}
   +C_{E,S,M,\chi}
      \left(\frac{A}{L}+\frac1L+\frac1A\right).
\end{equation}
\end{proposition}

\begin{equation}
   A\geq2\norm E_{C_tL^\infty},
   \qquad
   \frac{C_{\mathrm{alg},M}\norm E_{\A^M_{\R^2,t}}}{A}\leq\frac12.
\end{equation}

\begin{equation}\label{eq:Euclidean-amplitude-Wiener-bounds}
 \norm{a_L}_{\A^M_{\R^2,t}}
 +\norm{a_{L,\ell}}_{\A^0_{\R^2,t}}
 +\norm{h a_{L,\ell}}_{\A^0_{\R^2,t}}
 \leq C_{E,M,\chi}A^{1/2}.
\end{equation}

\begin{equation}\label{eq:Euclidean-frequency-sum-bound}
   \sum_{s\in\mu\Z\setminus\{0\}}
      \left(\sum_n\abs{c_nc_{s-n}}\right)\abs{s}^{1-S}
   \leq\mu^{1-S}\sum_{j\in\Z\setminus\{0\}}\abs j^{1-S}
   \leq C_S\mu^{1-S}.
\end{equation}

\begin{proposition}
\label{prop:Euclidean-mixed-absolute}
Fix $0<\eps<1$ and $S>2$.  Let $u\in C_tL^1(\R^2)$ satisfy
\[
   \norm{u}_{\A^0_{\R^2,t}}<\infty,
   \qquad
   \operatorname{AF}^{\partial_x}_S(u)<\infty,
   \qquad
   \widehat u(t,\xi,\eta)=0\quad\text{whenever }\abs\xi>K
\]
for every $t\in[0,1]$ and some $K\geq1$.  Under the hypotheses of
Proposition~\ref{prop:Euclidean-amplitude-exact-coefficient}, let
\[
   w=h a_{L,\ell}\rho_{\la,\eps},
\]
where $\rho_{\la,\eps}$ is given by Lemma~\ref{lem:slab}, and suppose
\begin{equation}\label{eq:Euclidean-mixed-separation}
   \mu\geq8\ell,
   \qquad
   \mu>8K.
\end{equation}
Then the projections of $\supp\widehat u$ and $\supp\widehat w$ onto the
$\xi$-axis are disjoint, and
\begin{equation}\label{eq:Euclidean-mixed-absolute}
\begin{split}
 \operatorname{AF}^{\partial_x}_S(u,w)
 &\leq C_{S,\eps}
   \norm{u}_{\A^0_{\R^2,t}}
   \norm{h a_{L,\ell}}_{\A^0_{\R^2,t}}
   \la^{-\frac12(1-\eps)}\mu^{1-S}\\
 &\leq C_{u,E,S,M,\eps,\chi}A^{1/2}
   \la^{-\frac12(1-\eps)-\eps(S-1)}.
\end{split}
\end{equation}
Consequently,
\begin{equation}\label{eq:Euclidean-combined-absolute}
\begin{split}
 \operatorname{AF}^{\partial_x}_S(u+w)
 \leq{}&\operatorname{AF}^{\partial_x}_S(u)
   +\norm{\partial_xE}_{\A^{-S}_{\R^2,t}}\\
 &+C_{E,S,M,\chi}\left(
       \frac A L+\frac1L+\frac1A+A\mu^{1-S}\right)\\
 &+C_{u,E,S,M,\eps,\chi}A^{1/2}
   \la^{-\frac12(1-\eps)-\eps(S-1)}.
\end{split}
\end{equation}
\end{proposition}

\begin{equation}\label{eq:Euclidean-mixed-w-Fourier}
 \widehat w(t,k)
 =\sum_{n\in\mu\Z\setminus\{0\}}c_n
    \widehat{h a_{L,\ell}}\bigl(t,k-(n,0)\bigr).
\end{equation}

\begin{proposition}\label{prop:Euclidean-onestep}
Fix $d\in\{3,5\}$, $\kappa\in\{-1,1\}$,
$0<\eps<1$, $0<\beta<\frac12(1-\eps)$, $S>d+1$, and $M>S+6$.
Let $u,E\in C^\infty([0,1];\mathcal S(\R^2))$ be a real pair satisfying
\eqref{eq:relaxed-general} with $r=2$.  Assume that $u$ and $E$ have compact spatial
Fourier support and that
\[
  \widehat u(t,\xi,\eta)=0\quad\text{for }\abs\xi<c_0
\]
for some $c_0>0$, and that
$\supp_tu\cup\supp_tE\subset I$ for a closed interval $I\Subset(0,1)$.
Let
\[
 K=\max\left(\{1,c_0\}\cup
 \{\abs\xi:(\xi,\eta)\in\supp\widehat u\}\right).
\]
Given finitely many exponents $1\leq p_j<2$ and numbers
$\delta_w,\delta_+>0$, one can choose parameters in the order
\[
 A\ \longrightarrow\ L\ \longrightarrow\ \ell\ \longrightarrow\ \la
\]
and a cutoff $h$ satisfying \eqref{eq:time-cutoff-lemma} with
$\tau=\la^{-\beta}$, for which $w=h a_{L,\ell}\rho_{\la,\eps}$ and
$u^+=u+w$, with $E^+$ defined by
\eqref{eq:Euclidean-update-general} for $r=2$, satisfy
\begin{align}
 \norm w_{C_tL^{p_j}(\R^2)}&<\delta_w,
 \label{eq:Euclidean-onestep-Lp}\\
 \norm{\partial_xE^+}_{\A^{-S}_{\R^2,t}}&<\delta_+,
 \notag\\
 \operatorname{AF}^{\partial_x}_S(u^+)
 &\leq\operatorname{AF}^{\partial_x}_S(u)
       +\norm{\partial_xE}_{\A^{-S}_{\R^2,t}}+\delta_w.
 \label{eq:Euclidean-onestep-absolute}
\end{align}
The pair $(u^+,E^+)$ is smooth and real, is Schwartz in space, has compact
spatial Fourier support, satisfies \eqref{eq:relaxed-general} with
$r=2$, and is
supported in the open $\la^{-\beta}$-neighborhood of $I$.  Furthermore,
\begin{equation}\label{eq:Euclidean-onestep-support}
 \supp\widehat w\subset
 \left\{(\xi,\eta):
   \mu-2\ell\leq\abs\xi\leq2\nu+2\ell,
   \ \abs\eta\leq2\ell\right\},
\end{equation}
where
\begin{equation}\label{eq:Euclidean-onestep-separation}
 \mu\geq8\ell,
 \qquad \mu>8K.
\end{equation}
Moreover,
\[
 \widehat{u^+}(t,\xi,\eta)=0\qquad(\abs\xi<c_0),
 \qquad
 \supp\widehat u\cap\supp\widehat w=\varnothing.
\]

Once $A,L,$ and $\ell$ are fixed, every sufficiently large admissible
$\la$ works.
\end{proposition}

\begin{proof}
Fix $(u,E)$, the finite list of exponents, and the two tolerances.
Choose $A\geq1$ so large that \eqref{eq:Euclidean-A-choice},
\eqref{eq:Euclidean-amplitude-smallness}, and
\begin{equation*}
 \frac{C_{E,S,M,\chi}}{A}<\frac{\delta_w}8
\end{equation*}
  hold.  Keeping this $A$ fixed, choose $L\geq1$ sufficiently large that
\begin{equation*}
 \frac{C_{E,\chi}+C_\chi A}{L}<\frac{\delta_+}{8},
 \qquad
 C_{E,S,M,\chi}\left(\frac{A}{L}+\frac1L\right)<\frac{\delta_w}8.
\end{equation*}
After $A$ and $L$ have been fixed, choose $\ell\geq1$ sufficiently large
that $C_{E,M,\chi}A\ell^{-M}<\delta_+/8$.

Choose an admissible $\la$ large enough that
\eqref{eq:Euclidean-onestep-separation} holds and
$\la^{-\beta}<\operatorname{dist}(I,\{0,1\})$.  Choose $h$ satisfying
\eqref{eq:time-cutoff-lemma} with $\tau=\la^{-\beta}$.  Then
\begin{equation*}
 0\leq h\leq1,\qquad h=1\ \text{on }I,\qquad
 \norm{h'}_{L^\infty}\leq C\la^\beta.
\end{equation*}
Since $E$ is supported in $I$, one has $hE=E$.  Define $w,u^+,$ and $E^+$
by \eqref{eq:Euclidean-increment} and
\eqref{eq:Euclidean-update-general}.

\subsubsection*{Dispersion error}

The dispersion error in $\partial_xE^+$ is
\[
 \partial_x\left((-1)^{(d+1)/2}\partial_x^{d-1}w
 +\kappa\partial_x^{-2}\partial_y^2w\right).
\]
By \eqref{eq:Euclidean-mixed-w-Fourier}, the Fourier support of $w$ is
contained in the disks
\[
 \{(\xi,\eta):\abs{(\xi-n,\eta)}\leq2\ell\},
 \qquad n\in\mu\Z\setminus\{0\},\quad \abs n\leq2\nu.
\]
On each disk, $\abs\xi\simeq\abs n$, $\abs\eta\leq2\ell$, and
$\langle(\xi,\eta)\rangle^{-S}\lesssim_S\abs n^{-S}$.
Since $\abs{\widehat w(t,\xi,\eta)}\leq\norm{w(t)}_{L^1}$, integration
over the disks and summation over $n\in\mu\Z\setminus\{0\}$ give
\begin{align*}
 \norm{\partial_x^dw}_{\A^{-S}_{\R^2,t}}
 &\leq C_S\norm w_{C_tL^1}
 \sum_{n\in\mu\Z\setminus\{0\}}
 \int_{\abs{(\xi-n,\eta)}\leq2\ell}
   \abs\xi^d\langle(\xi,\eta)\rangle^{-S}\,d\xi\,d\eta\\
 &\leq C_S\ell^2\norm w_{C_tL^1}
   \sum_{n\in\mu\Z\setminus\{0\}}\abs n^{d-S}
 \leq C_S\ell^2\mu^{d-S}\norm w_{C_tL^1},\\
 \norm{\partial_x(\partial_x^{-2}\partial_y^2w)}
      _{\A^{-S}_{\R^2,t}}
 &\leq C_S\norm w_{C_tL^1}
 \sum_{n\in\mu\Z\setminus\{0\}}
 \int_{\abs{(\xi-n,\eta)}\leq2\ell}
   \frac{\eta^2}{\abs\xi}\langle(\xi,\eta)\rangle^{-S}
   \,d\xi\,d\eta\\
 &\leq C_S\ell^4\norm w_{C_tL^1}
   \sum_{n\in\mu\Z\setminus\{0\}}\abs n^{-S-1}
 \leq C_S\ell^4\mu^{-S-1}\norm w_{C_tL^1}.
\end{align*}
Consequently,
\begin{equation}\label{eq:Euclidean-onestep-dispersion}
 \norm{\partial_x\bigl((-1)^{(d+1)/2}\partial_x^{d-1}w
       +\kappa\partial_x^{-2}\partial_y^2w)}_{\A^{-S}_{\R^2,t}}
 \leq C_S\left(\ell^2\mu^{d-S}
             +\ell^4\mu^{-S-1}\right)\norm w_{C_tL^1}.
\end{equation}
By \eqref{eq:Euclidean-Lp} with $p=1$, this is at most
\[
 C_{E,A,L,\ell,S,\eps}\left(
 \la^{\eps(d-S)-\frac12(1-\eps)}
 +\la^{-\eps(S+1)-\frac12(1-\eps)}\right),
\]
which tends to zero because $S>d+1$.

\subsubsection*{Nash error}

The Nash error is $\partial_x(2uw)$.  Minkowski's inequality and
Proposition~\ref{prop:Euclidean-mixed-absolute}, applied using
\eqref{eq:Euclidean-onestep-separation}, give
\begin{equation}\label{eq:Euclidean-onestep-Nash}
 \norm{\partial_x(2uw)}_{\A^{-S}_{\R^2,t}}
 \leq C_{u,E,S,M,\chi,\eps}A^{1/2}
 \la^{-\frac12(1-\eps)-\eps(S-1)}.
\end{equation}

\subsubsection*{Oscillation error}

The oscillation error is
\[
 \partial_x(E+w^2)
 =\partial_x\left[h^2\left((1-\chi_L^2)E+A\chi_L^2
 +(a_{L,\ell}^2-a_L^2)
 +a_{L,\ell}^2(\rho_{\la,\eps}^2-1)\right)\right].
\]
Put
$K_L=\widehat{\chi_L}*\widehat{\chi_L}$.  Since
$K_L\geq0$ and $\int_{\R^2}K_L=1$,
\[
 \widehat{(1-\chi_L^2)E}(t,k)
 =\int_{\R^2}K_L(z)
   \bigl(\widehat E(t,k)-\widehat E(t,k-z)\bigr)\,dz.
\]
Minkowski's inequality and the fundamental theorem of calculus give
\begin{align*}
 \int_{\R^2}
 \norm{\widehat{(1-\chi_L^2)E}(\cdot,k)}_{C_t}\,dk
 &\leq\int_{\R^2}\abs z
 (\widehat{\chi_L}*\widehat{\chi_L})(z)\,dz
 \int_0^1\int_{\R^2}
 \norm{\nabla\widehat E(\cdot,k-\sigma z)}_{C_t}\,dk\,d\sigma\\
 &\leq\frac{C_\chi}{L}
 \int_{\R^2}\norm{\nabla\widehat E(\cdot,k)}_{C_t}\,dk.
\end{align*}
Here the last step is \eqref{eq:Euclidean-spatial-cutoff-moment}.
Therefore,
\begin{equation*}
 \norm{\partial_x\bigl(h^2(1-\chi_L^2)E\bigr)}
      _{\A^{-S}_{\R^2,t}}
 \leq\frac{C_{E,\chi}}{L}.
\end{equation*}
The localized constant term satisfies
\begin{equation*}
 \norm{\partial_x(h^2A\chi_L^2)}_{\A^{-S}_{\R^2,t}}
 \leq\frac{C_\chi A}{L}
\end{equation*}
by \eqref{eq:localized-constant-scaling}.  Since $m_\ell=1$ on
$\{\abs k\leq\ell\}$, \eqref{eq:Euclidean-amplitude-Wiener-bounds}
gives
\[
 \norm{a_L-a_{L,\ell}}_{\A^0_{\R^2,t}}
 \leq\int_{\abs k>\ell}\norm{\widehat a_L(\cdot,k)}_{C_t}\,dk
 \leq\ell^{-M}\norm{a_L}_{\A^M_{\R^2,t}}
 \leq C_{E,M,\chi}A^{1/2}\ell^{-M}.
\]
The Wiener algebra inequality therefore yields
\begin{equation*}
 \norm{\partial_x\bigl(h^2(a_{L,\ell}^2-a_L^2)\bigr)}
      _{\A^{-S}_{\R^2,t}}
 \leq C_S\norm{a_{L,\ell}-a_L}_{\A^0_{\R^2,t}}
 \left(\norm{a_{L,\ell}}_{\A^0_{\R^2,t}}
       +\norm{a_L}_{\A^0_{\R^2,t}}\right)
 \leq C_{E,M,\chi}A\ell^{-M}.
\end{equation*}

For the remaining term, write
\[
 \rho_{\la,\eps}^2-1
 =\sum_{s\in\mu\Z\setminus\{0\}}
 \widehat{\rho_{\la,\eps}^2}(s)e^{2\pi isx}.
\]
The normalization gives $\widehat{\rho_{\la,\eps}^2}(0)=1$.  Since
\[
 \abs{\widehat{\rho_{\la,\eps}^2}(s)}
 \leq\sum_n\abs{c_nc_{s-n}},
\]
\eqref{eq:Euclidean-frequency-sum-bound} yields
\begin{equation}\label{eq:Euclidean-onestep-profile-convolution}
 \sum_{s\in\mu\Z\setminus\{0\}}
 \abs{\widehat{\rho_{\la,\eps}^2}(s)}\abs s^{1-S}
 \leq C_S\mu^{1-S}.
\end{equation}
Since $\supp\widehat{a_{L,\ell}^2}\subset\{\abs k\leq4\ell\}$ and
$\mu\geq8\ell$, for $s\neq0$ and $\abs k\leq4\ell$ one has
\[
 2\pi\abs{s+k_x}\left\langle(s,0)+k\right\rangle^{-S}
 \leq C_S\abs s^{1-S}.
\]
Therefore, \eqref{eq:Euclidean-onestep-profile-convolution} and
\eqref{eq:Euclidean-amplitude-Wiener-bounds} give
\begin{align*}
 \norm{\partial_x\left[
 h^2a_{L,\ell}^2(\rho_{\la,\eps}^2-1)\right]}
 _{\A^{-S}_{\R^2,t}}
 &\leq C_S\sum_{s\in\mu\Z\setminus\{0\}}
 \abs{\widehat{\rho_{\la,\eps}^2}(s)}\abs s^{1-S}
 \norm{a_{L,\ell}^2}_{\A^0_{\R^2,t}}\\
 &\leq C_{E,S,M,\chi}A\mu^{1-S}.
\end{align*}
These estimates and \eqref{eq:Euclidean-cancellation} yield
\begin{equation}\label{eq:Euclidean-onestep-oscillation}
 \norm{\partial_x(E+w^2)}_{\A^{-S}_{\R^2,t}}
 \leq\frac{C_{E,\chi}}{L}+\frac{C_\chi A}{L}
 +C_{E,M,\chi}A\ell^{-M}+C_{E,S,M,\chi}A\mu^{1-S}.
\end{equation}

\subsubsection*{Temporal error}

The temporal error in $\partial_xE^+$ is $\partial_tw$, where
\[
 \partial_tw=h'a_{L,\ell}\rho_{\la,\eps}
 +h(\partial_ta_{L,\ell})\rho_{\la,\eps}.
\]
Since $A-E\geq A/2$,
\[
 \partial_ta_L=-\frac{\chi_L\partial_tE}{2(A-E)^{1/2}},
 \qquad
 \partial_ta_{L,\ell}=P_{\leq\ell}^{x,y}\partial_ta_L,
\]
and hence
\[
 \norm{a_{L,\ell}}_{C_tW^{3,1}}
 +\norm{\partial_ta_{L,\ell}}_{C_tW^{3,1}}
 \leq C_{E,A,L,\ell}.
\]
Applying \eqref{eq:Euclidean-periodic-decoupling} with $p=1$ gives
\begin{align*}
 \norm{\partial_tw}_{C_tL^1}
 &\leq C\left(\norm{h'}_{L^\infty}
       \norm{a_{L,\ell}}_{C_tW^{3,1}}
       +\norm{\partial_ta_{L,\ell}}_{C_tW^{3,1}}\right)
       \norm{\rho_{\la,\eps}}_{L^1(\T^2)}\\
 &\leq C_{E,A,L,\ell,\eps}
       (1+\la^\beta)\la^{-\frac12(1-\eps)}.
\end{align*}
The gap in \eqref{eq:Euclidean-cone} defines $\partial_x^{-1}w$ and gives
$\partial_x\partial_t\partial_x^{-1}w=\partial_tw$.  Since $S>2$,
\begin{equation}\label{eq:Euclidean-onestep-temporal}
 \norm{\partial_x(\partial_t\partial_x^{-1}w)}
      _{\A^{-S}_{\R^2,t}}
 \leq\left(\int_{\R^2}\langle k\rangle^{-S}\,dk\right)
       \norm{\partial_tw}_{C_tL^1}
 \leq C_{E,A,L,\ell,S,\eps}
       (1+\la^\beta)\la^{-\frac12(1-\eps)}
 \longrightarrow0.
\end{equation}

By \eqref{eq:Euclidean-update-general},
\begin{align*}
 \norm{\partial_xE^+}_{\A^{-S}_{\R^2,t}}
 \leq{}&\norm{\partial_x(E+w^2)}_{\A^{-S}_{\R^2,t}}
 +\norm{\partial_x(2uw)}_{\A^{-S}_{\R^2,t}}\\
 &+\norm{\partial_x\bigl((-1)^{(d+1)/2}\partial_x^{d-1}w+
       \kappa\partial_x^{-2}\partial_y^2w)}_{\A^{-S}_{\R^2,t}}
 +\norm{\partial_x(\partial_t\partial_x^{-1}w)}
       _{\A^{-S}_{\R^2,t}}.
\end{align*}
For the absolute Fourier bound, Proposition~\ref{prop:Euclidean-mixed-absolute}
gives
\begin{align*}
 \operatorname{AF}^{\partial_x}_S(u^+)
 \leq{}&\operatorname{AF}^{\partial_x}_S(u)
 +\norm{\partial_xE}_{\A^{-S}_{\R^2,t}}\\
 &+C_{E,S,M,\chi}\left(
   \frac{A}{L}+\frac1L+\frac1A+A\mu^{1-S}\right)\\
 &+C_{u,E,S,M,\chi,\eps}A^{1/2}
 \la^{-\frac12(1-\eps)-\eps(S-1)}.
\end{align*}
The choices of $A,L,$ and $\ell$ control all $\la$-independent remainder
terms.  It remains to impose the separation and time-support conditions,
the finitely many bounds \eqref{eq:Euclidean-onestep-Lp}, the
$\la$-dependent terms in
\eqref{eq:Euclidean-onestep-dispersion}--
\eqref{eq:Euclidean-onestep-temporal}, and the two terms containing
$\mu$ or $\la$ in the last display.  The support conditions hold for all
sufficiently large admissible $\la$, and every listed norm or error term
tends to zero.  One sufficiently large choice of $\la$ therefore proves
\eqref{eq:Euclidean-onestep-Lp}--
\eqref{eq:Euclidean-onestep-absolute}.

\end{proof}

\end{document}
